\section{Bayes' Theorem}

If both $ P(A) $ and $ P(B) $ are nonzero, then from the definition of conditional probability, we can deduce that 
\begin{align}
	P(A \cap B) &=P(A)\cdot P(B|A)\\
	P(B \cap A) &=P(B)\cdot P(A|B),
\end{align}
and since $ P(A\cap B)= P(B\cap A) $, we conclude 
\begin{align}
	P(A)\cdot P(B|A) = P(B)\cdot P(A|B)
\end{align}
and therefore 
\begin{align}
	P(A|B)=\dfrac{P(A)\cdot P(B|A)}{P(B)}.
\end{align}
This result is known as \emph{Bayes' Theorem.}

\subsection{Generalizations}

The previous result can be extended in the following manner: Since 
\begin{align}
	\begin{cases}
		(B\cap A) \cup (B \cap A') &= B\\
		(B\cap A) \cap (B \cap A') &= \emptyset\\
	\end{cases},
\end{align}
then
\begin{align}
	P(B) & = P\left( (B\cap A) \cup (B \cap A') \right)\\
	&= P(B\cap A) + P(B\cap A')\\
	&= P(B|A)P(A) + P(B|A')P(A'),
\end{align}
so that 
\begin{align}
	P(A|B)=\dfrac{P(A)\cdot P(B|A)}{P(B|A)P(A) + P(B|A')P(A')}.
\end{align}

We can generalize this concept using a partition $ \seq{E_i}{i=0}{N} $ of the sample space $ S \supset A, B$. In that case
\begin{align}
	A = \bigcup_{i=0}^{N}(A\cap E_i)
\end{align}
and therefore
\begin{align}
	P(B) &= \sum_{i=0}^{N} P(B\cap E_i)\\
	&= \sum_{i=0}^{N} P(B|E_i)P(E_i).
\end{align}
We call this result the \emph{law of total probability} (or \emph{chain rule}).

Now, similarly to the first case of Bayes' theorem, we conclude that 
\begin{align}
	P(B|E_j)P(E_j)= P(E_j|B)P(B),
\end{align}
so using the law of total probability, we obtain
\begin{align}
	P(E_j|B)&= \dfrac{P(B|E_j)P(E_j)}{P(B)}\\
	&=\dfrac{P(B|E_j)P(E_j)}{\sum_{i=0}^{N} P(B|E_i)P(E_i)}
\end{align}

\begin{example}
	A gelatin manufacturing plant has three packaging machines. The distribution of packaging volume is carried out as follows:
	\begin{itemize}
		\item Machine 1: 38\%
		\item Machine 2: 32\%
		\item Machine 3: 30\%
	\end{itemize}
	
	Furthermore, the probability that a package comes out defective is 11\%, 15\%, and 14\% for each machine, respectively. 
	
	The plant's production management is interested in knowing the probability that if a unit is selected at random and found to be defective, it was packaged by Machine 2. 

	Let $ M_i $ denote the event that a gelatin unit was packaged by the $ i $-th machine, while $ D $ is the event that the unit is defective. 
	
	Calculate the probability that a gelatin unit came from the second machine given that it is defective. 
\end{example}
	
\begin{solution}
	
	According to the problem:
	\begin{align}
		P(M_1)&=0.38\\
		P(M_2)&=0.32\\
		P(M_3)&=0.30\\
		P(D|M_1)&=0.11\\
		P(D|M_2)&=0.15\\
		P(D|M_3)&=0.14
	\end{align}
	
	Then 
	\begin{align}
		P(M_2|D) &= \dfrac{P(D|M_2)P(M_2)}{	P(D|M_1)P(M_1)+P(D|M_2)P(M_2)+P(D|M_3)P(M_3)
		}\\
		&= \dfrac{(0.15)(0.32)}{
			(0.11)(0.38)+(0.15)(0.32)+(0.14)(0.30)	
		}\\
		&=0.3642=36.42\%
	\end{align}
\end{solution}
